\Needspace{13\baselineskip}
\chapterrule
\section{Attention：从直接读取到 Q/K/V 分工}
\label{intro:21c-naive-attention}
\label{intro:21b-why-attention}
\phantomsection\label{concept:attention}

\begin{idea}
\noindent\emph{按相关性读取上下文}\par

Attention（注意力）是一种按相关性读取信息的方法：
当前位置先对\textbf{可读取的每个位置}算一个匹配分数，
再把这些分数变成\textbf{权重}，
最后把各位置提供的信息按权重\textbf{加起来}。

它最核心的三个动作就是
\textbf{比较 $\to$ 分配权重 $\to$ 加权读取}。
\end{idea}

\subsection{从逐步传递到直接读取}

\conceptlink{rnn}{RNN} 让模型按顺序读，前文要经过一份份状态才能传到当前位置。
先比较两条信息路径：
% Side-by-side panels keep labels readable while fitting below the opening text.
\par\noindent\begin{minipage}{\linewidth}
\tagstructbegin{tag=Figure,alttext={左图 RNN 依次从春天到花再到开了传递状态；右图 Attention 从三个可见位置直接读取到位置三。右图线宽示意不同读取权重。}}
\begin{tcolorbox}[enhanced,colback=paper,colframe=linegray,
 boxrule=.7pt,arc=2mm,left=3mm,right=3mm,top=2mm,bottom=2mm]
\centering
\begin{tikzpicture}[>=Latex,x=1mm,y=1mm,
 word/.style={draw=linegray,fill=white,rounded corners=2pt,
  minimum width=17mm,minimum height=9mm,align=center,font=\footnotesize},
 output/.style={draw=flowcolor,fill=paper,rounded corners=2pt,
  minimum width=22mm,minimum height=10mm,align=center,font=\footnotesize}]
\node[font=\sffamily\bfseries\small,anchor=west] at (0,25) {RNN：依次传递};
\node[word] (r1) at (9,12) {春天\\$h_1$};
\node[word] (r2) at (30,12) {花\\$h_2$};
\node[word] (r3) at (51,12) {开了\\$h_3$};
\draw[->,very thick,flowcolor] (r1)--(r2);
\draw[->,very thick,flowcolor] (r2)--(r3);
\node[font=\footnotesize,align=center] at (30,-2) {较早信息经过中间状态};
\draw[linegray] (68,28)--(68,-9);
\node[font=\sffamily\bfseries\small,anchor=west] at (75,25) {Attention：直接读取};
\node[word] (a1) at (84,13) {春天\\$k_1,v_1$};
\node[word] (a2) at (108,13) {花\\$k_2,v_2$};
\node[word] (a3) at (132,13) {开了\\$k_3,v_3$};
\node[output] (at) at (108,-4) {位置 3 的\\读取结果};
\draw[->,flowcolor,line width=2pt] (a1.south) to[out=-80,in=165] (at.west);
\draw[->,flowcolor,line width=.6pt] (a2.south)--(at.north);
\draw[->,flowcolor,line width=1.15pt] (a3.south) to[out=-100,in=15] (at.east);
\end{tikzpicture}
\par
\end{tcolorbox}
\begin{center}\captionof{figure}{RNN 逐步传递与 Attention 直接读取的对照；线宽示意读取权重}\label{fig:attention-rnn}\end{center}
\tagstructend
\end{minipage}\par

RNN 的第 3 步能读到\textbf{上一份状态} $h_2$；
想重点利用第 1 个词的信息，就要靠这份状态此前保留了相关内容。

换一种设计：读到第三个词“开了”时，
让第三个位置直接查看位置 1、2、3，各自决定读多少。
图中右侧的线宽表示读取比例不同。接下来要解决的，就是怎样算出这些比例。

\subsection{从平均读取到按需读取}

\localfixture{二维输入，每行一个位置；与主线三维 Embedding 无关，本章后续投影对照与下一章批量计算继续沿用}
设当前位置可以读取三个位置，它们各自提供一条\textbf{内容向量}（行向量）：
\[
\vect v_1^\top=(1,0),\qquad
\vect v_2^\top=(0,1),\qquad
\vect v_3^\top=(1,1).
\]
最简单的方法是把它们\textbf{平均}：
\[
\vect y_{\rm avg}^\top
=\tfrac13\bigl(\vect v_1^\top+\vect v_2^\top+\vect v_3^\top\bigr)
=(2/3,\ 2/3).
\]
无论当前位置想查找什么，只要内容不变，平均结果就不变。

改为给三条内容分配权重 $a_1,a_2,a_3$，
并规定权重非负、总和为 $1$：
\begin{equation}
\boxed{\vect y^\top=a_1\vect v_1^\top+a_2\vect v_2^\top+a_3\vect v_3^\top,\qquad
a_i\ge0,\quad \sum_{i=1}^3 a_i=1.}
\label{eq:attention-design-goal}
\end{equation}
平均对应 $a_1=a_2=a_3=1/3$。
要让读取结果随当前位置的需要而变化，权重就得取决于
\textbf{当前想找什么}与\textbf{各位置凭什么被找到}的匹配程度。
于是每个位置在读取中承担三种职责：发起比较、被比较、交出内容。
本章先让同一条输入向量兼任这三种职责。

\subsection{三个二维输入，零参数读取}

把上一节的三条内容向量直接当作输入：
$\vect x_1^\top=(1,0)$，$\vect x_2^\top=(0,1)$，$\vect x_3^\top=(1,1)$。
每条输入既拿去比较，也被比较，还作为读出的内容，
也就是暂时令 $Q=K=V=X$，三种职责共用输入，不设置可训练投影。
这是用来观察读取过程的受限特例；后面会沿用同一组输入，把三种职责分开。
先将所有点积写成矩阵：
\[
X=\begin{bmatrix}1&0\\0&1\\1&1\end{bmatrix}
\ (3\times2),\quad
\underbrace{S}_{3\times3}=XX^\top
=\boxed{\begin{bmatrix}1&0&1\\0&1&1\\1&1&2\end{bmatrix}}.
\]
此处\textbf{未缩放、未加因果掩码}。第 $i$ 行是位置 $i$ 与三个位置的点积。
例如第三行 $(1,1,2)$：只对这条 Query 的三个分数计算 Softmax，得
$(0.212,0.212,0.576)$，继续把输入作为内容向量，
读出约为 $(0.788,0.788)$。
与平均读取的 $(0.667,0.667)$ 相比，结果偏向了匹配分数最高的位置 3。
不增加任何参数，也能完成一次“比较、分配、读取”。

\begin{boundarybox}
\textbf{零参数点积的匹配规则是固定的：}
(1) \textbf{分数对称。}$S=XX^\top=S^\top$：
“位置 $i$ 拿自己去比 $j$”与“位置 $j$ 拿自己去比 $i$”得到同一个原始分数。
(2) \textbf{长度也计入分数。}点积同时受方向和长度影响。
$\vect x_3^\top=(1,1)$ 最长，在三行里都拿到最高或并列最高的分数：
第一行 $(1,0,1)$、第二行 $(0,1,1)$、第三行 $(1,1,2)$。
逐行做 Softmax 后，第一行对位置 $1,3$ 给出并列最大权重；
第二行对位置 $2,3$ 给出并列最大权重；只有第三行对位置 $3$ 给出唯一最大权重。
这些分配由当前输入的方向、长度与固定的点积规则共同决定，尚未通过训练按任务需要调整。
两点同出一源：这一版本没有可调的数，训练无法让模型学会“按哪些特征匹配”。
\end{boundarybox}

要放宽这些限制，就让“拿什么询问”“凭什么被找到”“读出什么内容”
各由一组可训练投影产生。
下面就把三种职责拆开，并沿用同一组输入比较读出变化。
章末再给出训练 Q/K 选择所需记录的实验入口。
